\documentclass{article}
\usepackage{amsmath,hyperref}
\begin{document}
\section*{What a receiver may release after a late order update}
\paragraph{Result.}
A source-side publication fence does not exclude a release that was already in transit. A read of the accepted route does not prevent such a release arriving just after the read. In the tested model, a single receiver event that closes admission for the retired generation and reports the accepted prefix is sufficient to bound what still has authority. It is not proof that a footprint is clear. The construction preserves movement already admitted on a disjoint prefix. It uses an additional receiver-side operation, not a stock VDA cancellation procedure.

\paragraph{Source boundary.}
The \href{https://www.vda.de/dam/jcr%3Af0c9c019-1506-4dee-998a-e92723fbf025/EN-VDA5050-V2_1_0_2025-01.pdf}{VDA 5050 2.1 recommendation}, Sections 6.6.2--6.6.3, separates the released base from the modifiable horizon. An ordinary same-order update cannot erase an accepted base. This follows in the selected pinned ROS-Industrial VDA 5050 receiver too: validation checks both the decision-point identifier and sequence, and acceptance drops unreleased node and edge state rather than deleting released movement. The selected master defers an update while its previous update is not yet confirmed. The native fixture executes those source slices with an in-memory state and legal graph inputs assumed. It does not run the graph validator, action-state handling, transport, execution, or a vehicle.

The probe exposes a receiver-information issue, not an instruction for racing nominal master updates. Two admissible receiver states can reflect the same accepted update number and the same last traversed node but have different released prefixes. The contents bound to that update number differ between the worlds. When an authoritative retained ledger binds that number to a unique immutable payload, a fresh acknowledgment and known source law may already establish the accepted bound. The number in isolation is not the proof. The shim shows this by checking retained released nodes, the rejection of a stitch to a prior decision point, late lower-update rejection, and duplicate no-op.

\paragraph{Model and rule.}
Fix a single order with an append-only live base; each scoped old command is covered by a trusted accepted-prefix bound and affects only its mapped footprint. Higher authorities and other routes cannot act in the reserved footprint during the test. The new entrant requires exclusion of a conflict that could first arise at old step $j$. There are no resets, hidden replacement issuers, or sequence reuse. Every still-possible retired command has a generation below a known floor $f$. The proposed receiving control atomically raises its durable floor to $f$ and returns the accepted prefix $p$ from the same serialized event. An older command that is accepted first stays in $p$. An older command that arrives afterward cannot extend $p$. At coincident timestamps this distinction depends on serialization, not clock equality. Neither stating a floor without reading $p$ nor reading an unfenced $p$ is enough.

If $p<j$, the receipt can supply the unchanged frozen receiver's scoped exclusion fact. If $p\ge j$, the receipt alone cannot certify exclusion through the protection interval. A proof that an already-authorized conflicting part completed without possible re-entry, or an alternative lawful plan, may still succeed. The receipt is not a physical-clearance fact. With the separate sustained clearance and grant facts delivered by $q$ and $g$, fence evidence delivered by $a$, and protected duration $d$, the earliest start is $\max(a,q,g)$ in the model. This only works when each fact's actual guarantee lasts through the end of the protected operation. A missing or late fact preserves the refusal; an observation from another trace is not a substitute.

\paragraph{Comparison and exclusions.}
The source check executes six client branches and one master guard. The model replay runs the byte-identical frozen receiver from \href{https://github.com/eternal-november/eos-research-programmes/pull/462}{PR 462}, and agrees with a direct serial check in 258,048 parameter combinations. In the named case it admits launch at tick 2 and completion at 4, rather than missing the deadline waiting until tick 5 for stipulated reconciliation. Acceptance of the conflicting old step first, or clearance delayed to tick 3, defeats that guarantee. The wait time and plant footprint are authored; these are not measured native performance or field counts. With the same effective atomic receipt, the reference ties. Resource-side generation checking is established (for example \href{https://static.usenix.org/events/osdi06/tech/full_papers/burrows/burrows_html/}{Burrows, Section 2.4}); this probe connects it to the fact needed by a live scheduling receiver rather than claiming a new locking protocol. It does not yet establish an advantage over an applicable published controller using the same resource operation, an independent transfer, or a way to synthesize the footprint map from source.

Before transferring the receipt, validate generation scope, durable floor, correlation of the reported $p$ to the authoritative accepted state, and no motion outside that state. Then check foundation, motion, and other near-bay occupants separately. In particular a reset that loses the floor, an older issued physical effect still in progress, or any later authorized entrant admitted to the same scope invalidates the claimed certificate. No stock command or live deployment is recommended.
\end{document}
